% !PRINTED 635
\weberpage{635}

In der Betrachtung dieses Abschnittes haben wir, dem Hauptziel der Untersuchung
entsprechend, die Koeffizienten $a_1, a_2, \ldots, a_n$ \gesperrt{als rationale Funktionen von} $z$
betrachtet. Alles bleibt aber vollst\"andig unge\"andert, wenn wir f\"ur diese Gr\"o\ss{}en
irgend einen \gesperrt{Rationalit\"atsbereich} festsetzen.

\begin{center}
\S\ 175. \gesperrt{Ganze Funktionen von $z$.}
\end{center}

Jede Funktion $\omega$ des K\"orpers $\Omega$ gen\"ugt, wie wir gesehen haben, einer
Gleichung niedrigsten Grades:

\weq{1}{\omega^e + b_1\omega^{e-1} + b_2\omega^{e-2} + \cdots + b_e = 0,}

deren Koeffizienten $b_1, b_2, \ldots, b_e$ rationale Funktionen von $z$ sind.

1.\ \gesperrt{Wenn diese Koeffizienten ganze rationale Funktionen von $z$ sind, so hei\ss{}t
$\omega$ eine ganze algebraische oder kurz eine ganze Funktion von $z$.}

\"Uber die ganzen algebraischen Funktionen gelten die n\"amlichen S\"atze, wie \"uber die
ganzen Zahlen (Bd.\ II, \S\,149).

Setzen wir

\weq{2}{\varphi(t) = t^e + b_1 t^{e-1} + b_2 t^{e-2} + \cdots + b_e,}

so ist, wie wir in \S\,172 gesehen haben,

\weq{3}{N(t - \omega) = \varphi(t)^f \quad (ef = n)}

eine ganze Potenz von $\varphi(t)$. Wenn wir daher $N(t-\omega)$ nach Potenzen von $t$
ordnen, so werden die Koeffizienten alle wieder ganze rationale Funktionen, also
insbesondere:

2.\ \gesperrt{Die Norm und die Spur einer ganzen Funktion $\omega$ sind ganze rationale
Funktionen von $z$.}

3.\ \gesperrt{Eine rationale Funktion von $z$ ist nur dann eine ganze algebraische
Funktion, wenn sie eine ganze rationale Funktion ist.}

Denn ist $\omega = -b$ rational, so ist $e = 1$ und $\omega + b = 0$ die Gleichung
niedrigsten Grades f\"ur $\omega$, also $\omega$ nur dann ganz, wenn $b$ ganz und rational
ist.

4.\ \gesperrt{Jede Funktion $\eta$ in $\Omega$ kann durch Multiplikation mit einer von Null
verschiedenen rationalen Funktion von $z$ in eine ganze algebraische Funktion verwandelt
werden.}

Denn jede Funktion $\eta$ gen\"ugt einer Gleichung niedrigsten Grades:

\weq{4}{b_0\eta^e + b_1\eta^{e-1} + \cdots + b_{e-1}\eta + b_e = 0,}
%CONT>>
